Combining formal verification strategies---grammar constraints, static analysis, and SMT-guided repair---promises multiplicative improvements in LLM-generated code correctness. But does this work in practice? We provide the first unified comparison of these strategies, measuring whether they target independent error classes. Our experiments on HumanEval reveal that error classes are largely independent (mean Jaccard index 0.0752), validating layered verification pipelines. Grammar-constrained decoding reduces syntax errors by 40\%. However, static analysis feedback loops fail when using completion models: a model given repair feedback \emph{increased} security issues by 600\%. This reveals a critical prerequisite: detection is model-agnostic, but repair requires instruction-tuned models. Our error-class-independence framework provides principled guidance for combining verification strategies, while our failure analysis identifies a capability requirement that practitioners must verify before deploying feedback-based pipelines.
